\documentclass[11pt,a4paper]{article}

% ---------- Packages ----------
\usepackage[margin=1in]{geometry}
\usepackage{amsmath, amssymb}
\usepackage{graphicx}
\usepackage{booktabs}
\usepackage{caption}
\usepackage[hidelinks]{hyperref}
\usepackage{siunitx}
\usepackage{float}
\usepackage{authblk}

% Tight spacing
\setlength{\parskip}{10pt}
\setlength{\parindent}{0pt}
\renewcommand{\baselinestretch}{1.15}
\captionsetup{font=small, labelfont=bf, skip=3pt}

% ---------- Shortcuts ----------
\newcommand{\dd}{\mathrm{d}}
\newcommand{\E}{\mathbb{E}}
\newcommand{\Var}{\mathrm{Var}}
\newcommand{\R}{\mathbb{R}}

% Tighter section headings
\usepackage{titlesec}
\titleformat{\section}{\normalsize\bfseries}{\thesection.}{0.5em}{}
\titlespacing*{\section}{0pt}{6pt}{2pt}

% ---------- Title ----------
\title{\vspace{-2.5em}\textbf{\large Pairs Trading via Ornstein--Uhlenbeck Mean Reversion}\\[-2pt]
{\normalsize A Continuous-Time Stochastic Calculus Approach}}
\author{\normalsize Vedika Jain- 2024B4A81111G, Meemansa Jha- 2024B4A71096G}
\affil{\footnotesize Stochastic Calculus for Finance -- Course Project}
\date{\vspace{-1.5em}}

\begin{document}
\maketitle
\vspace{-2.8em}

% ==========================================================
\section{Motivation and Theoretical Framework}
% ==========================================================

Individual equity prices are typically modelled as geometric Brownian motions
and are therefore non-stationary. Yet for certain pairs of stocks in the same
industry, a linear combination of their log-prices is empirically
\emph{stationary}. Such pairs are \emph{cointegrated} (Engle \& Granger, 1987);
the stationary combination (the \emph{spread}) is a natural candidate for a
mean-reverting SDE. We model the spread as an Ornstein--Uhlenbeck (OU) process,
calibrate it to training data, and trade a rule that enters when the spread
is far from its mean and exits when it reverts.

\paragraph{The OU SDE.} On $(\Omega, \mathcal{F}, \{\mathcal{F}_t\}, \mathbb{P})$
supporting a standard Brownian motion $W(t)$ (Shreve~\cite{shreve2004}, Ch.~3), the
OU process $X(t)$ satisfies
\begin{equation}
    \dd X(t) = \theta\bigl(\mu - X(t)\bigr)\,\dd t + \sigma\,\dd W(t), \qquad X(0)=x_0,
    \label{eq:ou}
\end{equation}
with $\theta>0$ (mean-reversion speed), $\mu\in\R$ (long-run mean), $\sigma>0$
(volatility). Equation~\eqref{eq:ou} is the Vasicek SDE of Shreve
Example~4.4.10 under $\theta\!=\!\beta,\mu\!=\!a/\beta$.

\paragraph{Closed-form solution via the It\^o--Doeblin formula.}
Let $Y(t)=e^{\theta t}X(t)$. Applying the It\^o--Doeblin formula
(Shreve Thm.~4.4.1) to $f(t,x)=e^{\theta t}x$ ($f_t=\theta e^{\theta t}x$,
$f_x=e^{\theta t}$, $f_{xx}=0$) and substituting \eqref{eq:ou}, the
$X(t)\,\dd t$ terms cancel to leave
$\dd Y(t) = \theta\mu e^{\theta t}\dd t + \sigma e^{\theta t}\dd W(t)$.
Integrating and dividing by $e^{\theta t}$ gives
\begin{equation}
    X(t) = x_0 e^{-\theta t} + \mu\bigl(1-e^{-\theta t}\bigr)
           + \sigma\!\int_0^t e^{-\theta(t-s)}\,\dd W(s).
    \label{eq:ou-sol}
\end{equation}
The first two terms describe deterministic exponential decay of the initial
deviation toward $\mu$; the third is an It\^o integral (Shreve~\S4.3) against
the exponential kernel. By the It\^o isometry
$\Var[X(t)]=(\sigma^2/2\theta)(1-e^{-2\theta t})$, so as $t\to\infty$,
$X(t)\Rightarrow\mathcal{N}(\mu,\sigma^2/(2\theta))$. The half-life of
mean reversion is $\tau=\ln 2/\theta$, and the stationary standard deviation
$\sigma_s=\sigma/\sqrt{2\theta}$ quantifies what ``far from the mean'' means.

\paragraph{From GBM prices to a stationary spread.} Each $\log S_i(t)$ has
unbounded variance. But if two log-prices share a common stochastic trend,
a cointegrating vector $(1,-\beta)$ exists such that
$Z(t)=\log S_A(t)-\alpha-\beta\log S_B(t)$ is stationary, and
$Z(t)$ can be modelled by \eqref{eq:ou}.

% ==========================================================
\section{Methodology}
% ==========================================================

\textbf{Data.} Daily adjusted-close prices for eight candidate US/Indian pairs
(Yahoo Finance, 2021-01-01 to 2024-12-31). First 75\% of days used for
estimation (\emph{training}), last 25\% reserved for \emph{out-of-sample}
backtest. No test data informs the screen, the fit, or the thresholds.

\textbf{Step 1 -- Cointegration screen.} For each pair we run the
Engle--Granger two-step procedure: OLS of $\log S_A$ on $\log S_B$ yields
$(\alpha,\beta)$; an Augmented Dickey--Fuller test is then run on the residuals.
Accept at the 5\% level if ADF $<-3.34$ (MacKinnon).

\textbf{Step 2 -- OU calibration.} The exact discretisation of \eqref{eq:ou}
is AR(1): $X_{t+\Delta t}=aX_t+b+\eta_t$ with
$a=e^{-\theta\Delta t},\;b=\mu(1-a)$. OLS of $X_{t+\Delta t}$ on $X_t$
with $\Delta t=1/252$ gives $\theta=-\ln(a)/\Delta t$, $\mu=b/(1-a)$,
$\sigma=\sqrt{2\theta\,\Var(\eta)/(1-a^2)}$.

\textbf{Step 3 -- Trading rule.} Let $Z_t=(X_t-\mu)/\sigma_s$. \textbf{Enter}
when $|Z_t|>2$ (short spread if $Z_t>2$, long if $Z_t<-2$);
\textbf{take profit} at $|Z_t|<0.5$; \textbf{stop-loss} at $|Z_t|>3.5$;
\textbf{time-stop} after $2\tau$. Re-entry is blocked until $|Z_t|$ re-enters
$\pm 2$ (prevents whipsaw). Thresholds are set \emph{a priori} from the
stationary distribution of $Z$ and never tuned.

\textbf{Costs.} Leg $A$ at unit notional, leg $B$ at $\beta$; 10 bp per leg per
trade, so round-trip cost $=2(1+|\beta|)\times10$ bp.

% ==========================================================
\section{Empirical Results}
% ==========================================================

\textbf{Pair screening.} Table~\ref{tab:screen} reports the eight pairs.
Two pass the 5\% cointegration test: \textbf{MA/V} ($\text{ADF}=-3.77$,
half-life 19.5 d) and \textbf{RELIANCE.NS/ONGC.NS} ($\text{ADF}=-3.39$,
half-life 25.7 d). The six rejected pairs confirm that sectoral similarity
does not imply cointegration.

\begin{table}[H]
\centering\footnotesize
\caption{Engle--Granger test and OU parameters (training window, 756 days). Bold rows pass at 5\%.}
\label{tab:screen}
\setlength{\tabcolsep}{5pt}
\begin{tabular}{lS[table-format=1.3]S[table-format=-1.2]cS[table-format=1.2]S[table-format=1.3]S[table-format=3.1]}
\toprule
Pair & {$\beta$} & {ADF} & Coint. & {$\theta$\,/yr} & {$\sigma_s$} & {Half-life (d)} \\
\midrule
KO/PEP                   & 0.738 & -2.46 & No  & 3.72 & 0.036 &  47.0 \\
XOM/CVX                  & 1.243 & -1.76 & No  & 1.92 & 0.082 &  91.1 \\
\textbf{MA/V}            & \textbf{0.952} & \textbf{-3.77} & \textbf{Yes} & \textbf{8.96} & \textbf{0.029} & \textbf{19.5} \\
GLD/SLV                  & 0.231 & -1.30 & No  & 1.57 & 0.056 & 111.6 \\
HD/LOW                   & 0.839 & -2.31 & No  & 3.85 & 0.046 &  45.3 \\
HDFCBANK.NS/ICICIBANK.NS & 0.248 & -3.15 & No  & 6.31 & 0.055 &  27.7 \\
TCS.NS/INFY.NS           & 0.467 & -2.02 & No  & 3.48 & 0.057 &  50.1 \\
\textbf{RELIANCE.NS/ONGC.NS} & \textbf{0.336} & \textbf{-3.39} & \textbf{Yes} & \textbf{6.79} & \textbf{0.061} & \textbf{25.7} \\
\bottomrule
\end{tabular}
\end{table}

\textbf{Backtest (out-of-sample, 252 days).} Table~\ref{tab:bt} gives the
aggregate results. Three trades: MA/V produced one clean reversion win
and one stop-loss; RELIANCE/ONGC produced one small win. Combined success
ratio $2/3=66.7\%$, net P\&L $+2.07\%$ after costs.

\begin{table}[H]
\centering\footnotesize
\caption{Out-of-sample backtest results (2023-11-27 to 2024-12-31).}
\label{tab:bt}
\setlength{\tabcolsep}{6pt}
\begin{tabular}{lccS[table-format=1.3]S[table-format=-1.2]S[table-format=-1.2]S[table-format=-1.2]S[table-format=-1.2]S[table-format=2.0]}
\toprule
Pair & \# trades & wins & {success} & {avg \%} & {total \%} & {Sharpe} & {max DD\%} & {avg hold (d)} \\
\midrule
MA/V                 & 2 & 1 & 0.500 & 1.00 &  2.01 & 0.53 & -2.67 & 33 \\
RELIANCE.NS/ONGC.NS  & 1 & 1 & 1.000 & 0.07 &  0.07 & {--} &  0.00 & 26 \\
\midrule
\textbf{Combined}    & \textbf{3} & \textbf{2} & \textbf{0.667} & \textbf{0.69} & \textbf{2.07} & {--} & \textbf{-2.67} & \textbf{31} \\
\bottomrule
\end{tabular}
\end{table}

\textbf{Illustration -- MA/V.} Figure~\ref{fig:main} (left) shows the test-set
spread and $z$-score with trade markers. Trade~1 (Mar--May 2024, green)
is a textbook reversion: entry at $z\!\approx\!+2.2$, the spread relaxes
to $\mu$ over $\sim\!30$ days, exit at $+0.5\sigma_s$ for $+4.7\%$.
Trade~2 (Aug--Sep 2024, red) is the opposite: after entry at
$z\!\approx\!+2.5$, the spread drifts to $+4\sigma_s$, triggering the
$3.5\sigma_s$ stop at $-2.7\%$. The top panel shows the spread breaking
through the stop band before eventually reverting in November, consistent
with a temporary cointegration breakdown -- precisely the scenario the
stop-loss is designed to contain. Figure~\ref{fig:main} (right) plots the
cumulative P\&L curves.

\begin{figure}[H]
\centering
\includegraphics[width=0.58\linewidth]{best_pair_spread.png}\hfill
\includegraphics[width=0.40\linewidth]{equity_curves.png}
\caption{\textit{Left:} MA/V out-of-sample spread (top) and $z$-score (bottom).
Dots = entries, crosses = exits. \textit{Right:} cumulative out-of-sample P\&L
by pair, net of costs.}
\label{fig:main}
\end{figure}

% ==========================================================
\section{Discussion and Limitations}
% ==========================================================

\textbf{What worked.} The cointegration filter has teeth: six of eight
sector-matched candidates were rejected, and both retained pairs were
profitable out-of-sample. The MA/V winner matches the OU prediction almost
exactly -- entry just above $+2\sigma_s$, monotone drift to $\mu$, take-profit
exit -- showing that when the stationarity assumption holds, the edge is real.

\textbf{What failed.} The MA/V stop-loss exposes the main model risk: OU
parameters estimated in training need not remain valid out-of-sample.
Cointegration can drift or break under regime shifts. The $3.5\sigma_s$ stop
contained the loss at $-2.7\%$; without it, the position would have stayed
open for roughly three additional months in drawdown.

\textbf{Key limitations.} \emph{Sample size:} three trades across two pairs
is too few for precise claims; $66.7\%$ is directional evidence, not a
point estimate. \emph{Survivorship bias:} candidates were chosen from
well-known sector analogues, biasing the screen toward relationships that
remain recognisable today. \emph{Static $\beta$:} a Kalman-filtered hedge
ratio that updates online would track drift and likely reduce stop-out
frequency. \emph{Cost sensitivity:} at 20 bp/leg (more representative of
Indian equities with STT) roughly half the MA/V profit would be eroded and
the RELIANCE trade would turn negative.

\textbf{Conclusion.} The OU model, derived directly from the
It\^o--Doeblin formula applied to the cointegration residual, provides a
coherent bridge from the continuous-time machinery in Shreve Chs.~3--4 to a
concrete trading strategy. The out-of-sample evidence is modestly positive
($+2.07\%$ net in one year, $2/3$ hit rate) and, more importantly, behaves
as theory predicts: one clean reversion, one stop-loss where the
stationarity assumption temporarily failed. Together these two outcomes
form a compact demonstration of both the power and the limits of
continuous-time mean-reversion modelling in practice.

% ==========================================================
\vspace{-2pt}
{\footnotesize
\begin{thebibliography}{9}
\setlength{\itemsep}{-1pt}
\bibitem{shreve2004} Shreve, S.\,E. (2004). \emph{Stochastic Calculus for Finance~II: Continuous-Time Models.} Springer.
\bibitem{eg1987} Engle, R.\,F.\ \& Granger, C.\,W.\,J. (1987). Co-integration and error correction. \emph{Econometrica} 55(2), 251--276.
\bibitem{gatev2006} Gatev, E., Goetzmann, W.\,N.\ \& Rouwenhorst, K.\,G. (2006). Pairs trading: performance of a relative-value arbitrage rule. \emph{RFS} 19(3), 797--827.
\bibitem{mackinnon2010} MacKinnon, J.\,G.\ (2010). Critical values for cointegration tests. \emph{Queen's Econ.\ WP} No.~1227.
\end{thebibliography}
}

\end{document}
